\documentclass[a4paper,12pt]{jreport}
\usepackage{amsmath}
\usepackage[dvipdfmx]{graphics}
\usepackage{lscape}

%\documentstyle[a4j,12pt,epsf]{jreport}

% margine control
\setlength{\textwidth}{16.5cm} \setlength{\textheight}{23.5cm}
\setlength{\oddsidemargin}{10pt} \setlength{\evensidemargin}{10pt}
\setlength{\topmargin}{-10pt} \setlength{\marginparsep}{15pt}
\setlength{\parskip}{0.3cm}

% reassignment of bib
\renewcommand{\bibname}{参考文献}




\begin{document}

\begin{titlepage}
  \centering

  % ===== 上段 =====
  \vspace*{22mm}
  {\Large 2025年度 秋学期\par}

  \vspace{1mm}
  {\Large 卒\ \ 業\ \ 論\ \ 文\par}

  % ===== タイトル =====
  \vspace{15mm}
  {\LARGE\bfseries
   生地拡大画像と組成情報テキストによる\\マルチモーダルモデルを用いた衣服の着心地推定
  \par}

  % ===== 指導教員 =====
  \vspace{22mm}
  {\Large 指導教員：島川 博光\par}

  % ===== 所属 ====
  {\Large 立命館大学 情報理工学部\par}
  \vspace{1mm}
  {\Large 卒業研究３（AC）\par}

  % ===== 下段の情報 =====
  \vspace{22mm}
  {\large コース：システムアーキテクト（4回生）\par}
  \vspace{1mm}
  {\large 学生証番号：2600220391-9\par}
  \vspace{1mm}
  {\large 氏名：森田\ \ 欣司\par}

  \vspace*{10mm}
\end{titlepage}




\pagenumbering{roman}

\begin{center}
{\LARGE {\bf 生地拡大画像と組成情報テキストによる\\マルチモーダルモデルを用いた衣服の着心地推定}}


\vspace{10mm}

{\Large {\bf 概要}}
\end{center}

\addcontentsline{toc}{chapter}{概要}
\setlength{\baselineskip}{8mm}

近年，インターネット上のショッピングサイトを利用した衣服購入が一般化しているが，実際に生地に触れることができないため，購入前に着心地を把握することは困難である．この問題に対し，本研究では，生地拡大画像および組成情報を用いて，衣服の着心地を推定する手法を提案する．

着心地は主観的かつ多面的な概念であり，明確な数値ラベルを付与することが難しい．そこで本研究では，レビュー文に含まれる着心地表現を教師信号として利用し，着心地を低次元の連続値ベクトルとして表現する枠組みを採用した．提案手法では，生地拡大画像から畳み込みニューラルネットワークを用いて視覚的特徴を抽出し，組成情報から多層パーセプトロンを用いて素材構成に基づく特徴を抽出する．これらの特徴をゲート機構を用いたマルチモーダル学習により統合し，衣服の着心地を表現するベクトルを推定する．

学習においては，レビュー文を文埋め込みモデルによってベクトル化し，これを教師データとして用いた．大規模なレビュー全文を用いた事前学習と，着心地に関する記述のみを抽出した高品質データによる再学習からなる二段階学習を行うことで，着心地表現の精度向上を図った．

実験の結果，再学習を行うことで，推定された着心地ベクトルと教師ベクトルとの平均コサイン類似度が向上することを確認した．また，近傍分析により，生地の重さや柔らかさなど，視覚的特徴や素材構成と関連の強い着心地表現については，意味的に類似した衣服がベクトル空間上で近接して配置されることを示した．以上の結果から，本研究で提案した手法が，生地拡大画像および組成情報から着心地を推定する上で有効であることが示唆された．

\clearpage


\begin{center}
{\LARGE {\bf Estimating Clothing Comfort with a Multimodal Model Based on Fabric Close-up Images and Material Composition Information}}


\vspace{8mm}

{\Large {\bf abstract}}
\end{center}

In online fashion shopping, customers cannot physically touch garments before purchase, which makes it difficult to assess their wearing comfort in advance. To address this issue, this study proposes a method for estimating clothing comfort using fabric close-up images and material composition information.

Comfort is a subjective and multi-faceted concept, and it is difficult to define explicit numerical labels for supervision. Therefore, this study represents clothing comfort as a low-dimensional continuous vector by leveraging comfort-related expressions contained in customer reviews as weak supervisory signals. In the proposed approach, visual features are extracted from fabric close-up images using a convolutional neural network, while material composition features are extracted from fabric composition data using a multilayer perceptron. These features are integrated through a gated multimodal learning framework to estimate a comfort representation vector for each garment.

For training, review texts are converted into vector representations using a sentence embedding model and used as teacher vectors. A two-stage learning strategy is employed, consisting of pretraining with a large-scale dataset constructed from full review texts and fine-tuning with a smaller, manually curated dataset focusing on comfort-related descriptions. The experimental result showed that fine-tuning improves the cosine similarity between the estimated comfort vectors and the teacher vectors. Furthermore, the nearest-neighbor analysis demonstrated that garments with semantically similar comfort characteristics, such as softness and fabric weight, are located close to each other in the embedding space. These results suggest that the proposed method is effective for estimating clothing comfort from visual and material information alone.
\addcontentsline{toc}{chapter}{abstract}


\clearpage


\addcontentsline{toc}{chapter}{目次}
\tableofcontents

\clearpage
\addcontentsline{toc}{chapter}{図目次}
\listoffigures
\clearpage
\addcontentsline{toc}{chapter}{表目次}
\listoftables

%%%%%%%%%%%%%%%%  End of Prologue   %%%%%%%%%%%%%%%%%%%%%%%%%%%%%%%%%%

\clearpage
\pagenumbering{arabic}

\chapter{はじめに}

近年，インターネット上のショッピングサイトを利用した衣服の購入が広く普及している．オンラインでの衣服購入では，実店舗と異なり，実際に衣服を手に取って触れることができないため，生地の質感や着心地を事前に把握することが難しい．その結果，購入後に着心地が想定と異なると感じる場合があり，これはオンラインショッピングにおける重要な課題の一つとなっている\cite{Park}．

このような課題に対し，商品画像や商品説明文レビューなど，オンライン上で提供される情報を活用し，衣服の特徴や使用感を推定する試みが行われてきた．特に，衣服画像を用いた研究では，衣服カテゴリ，素材，質感といった視覚的属性を推定する手法が提案されている\cite{LiuZ}\cite{Bell}．これらの研究では，主に畳み込みニューラルネットワーク(CNN)を用いて，衣服画像から特徴を抽出することで，衣服の外観に関する情報を数値的に表現している．

一方，レビューを用いた研究では，ユーザが記述したテキストから，衣服に対する印象や主観的評価を推定する試みが行われている\cite{McAuley15}\cite{He}．これらの研究は，衣服に対する感性や嗜好を扱う上で有効であるが，多くの場合，推論時にもレビューを必要とするという制約を持つ．

さらに，画像やテキストといった異なる形式の情報を統合する，
マルチモーダル学習を用いた衣服解析も注目されている．マルチモーダル学習は，単一の情報源のみを用いる場合と比較して，より多面的な衣服表現を獲得できる点で有効であると報告されている\cite{Ngiam}\cite{Baltrušaitis}\cite{Ni}．

一方で，衣服の着心地は，生地の柔らかさや厚み，伸縮性といった物理的特性に加え，使用環境や個人差にも影響される多面的な概念である．このため，着心地を単一の指標として定義し，自動的に推定することは容易ではない\cite{Norman}．そのため，実利用に即した情報を活用しつつ，着心地を計算機上で扱うための表現方法について検討することが求められている．

近年では，深層学習を用いたファッション推薦や衣服解析に関する研究が進展しており，衣服画像，レビュー文，属性情報などを統合的に扱う手法が多数提案されている．これらの研究では，画像とテキストを組み合わせることで，衣服の印象や好みをより精緻に捉えることが可能であることが示されている\cite{Ni}\cite{McAuley13}\cite{Zheng}．また，2020年代に入ってからは，深層学習を用いたファッション推薦手法を体系的に整理したレビュー研究も報告されており，マルチモーダル学習が重要な役割を果たしていることが指摘されている \cite{LiuQ}\cite{Wu}．しかし，これらの研究の多くは，推薦や分類タスクを主目的としており，推論時にもレビューやテキスト情報を必要とする場合が多い．そのため，画像や組成といった商品情報のみから，着心地そのものを低次元の連続表現として推定する枠組みについては，十分に検討されていない．

本研究の目的は，生地拡大画像および組成情報を用いて，衣服の着心地を推定する手法を構築することである．特に，レビューに含まれる着心地表現を教師信号として利用し，視覚情報と数値情報を統合することで，着心地をベクトルとして表現する枠組みを検討する．

提案手法では，生地拡大画像から得られる視覚的特徴と，組成情報から得られる数値的特徴をマルチモーダル学習によって統合し，衣服の着心地を低次元の連続値ベクトルとして表現することを目指す．まず，生地拡大画像および組成情報をそれぞれ独立した特徴抽出器によりベクトル化する．得られた特徴を統合層に入力することで，着心地を表す共通表現を学習する．これにより，異なる衣服間の着心地の類似性を定量的に扱うことを可能とする．実験の結果，近傍分析により，同様の着心地を有する衣服がベクトル空間上で近接して配置されることを確認した．

本論文の構成を以下に示す．
第2章では，本研究で用いるCNN，マルチモーダル学習，およびテキスト埋め込みといった基礎技術について説明する．
第3章では，生地拡大画像および組成情報を用いた着心地推定手法の全体構成について述べる．
第4章では，本研究で使用したデータセットの概要およびデータ処理の流れについて説明する．
第5章では，実験設定について述べる．
第6章では，実験結果を示し，提案手法の有効性および限界について考察する．
最後に第7章では，本研究のまとめと今後の課題について述べる．



\chapter{着心地推定のためのマルチモーダル学習}

\section{衣服画像を用いた属性推定}

衣服画像を用いた属性推定に関する研究では，画像中に含まれる形状や質感，色合いなどの視覚的特徴を抽出するために，CNNが広く用いられている\cite{LeCun}．CNNは，局所的なパターンを段階的に学習する構造を持ち，画像分類や物体認識などの分野において高い性能を示している．

衣服画像に対する属性推定では，カテゴリ分類や色，柄，素材といった属性の推定が行われており，CNNを用いることでこれらの特徴を自動的に抽出することが可能である．特に，生地の凹凸や織り方などの微細な視覚情報は，画像全体を一様に扱う手法よりも，局所的な特徴を捉えるCNNによって効果的に表現できる．衣服画像を用いた解析では，CNNによる特徴抽出が広く用いられており，衣服のカテゴリや質感を捉える手法が提案されている\cite{LiuZ}\cite{Bell}\cite{Cimpoi}．生地拡大画像に含まれる質感や構造的特徴は，人間が知覚する素材感とも関連していることが報告されている\cite{Sharan}．

このような理由から，衣服画像から特徴を抽出する手法として，CNNが適している．CNNを用いることで，生地拡大画像に含まれる質感や構造的特徴を数値ベクトルとして表現することが可能となる．

\section{マルチモーダル学習による衣服解析}
衣服に関する情報は，画像や数値データなど，異なる形式の情報として与えられることが多い．これらの異種情報を同時に扱うための枠組みとして，マルチモーダル学習が用いられている\cite{Ngiam}．マルチモーダル学習では，複数の情報源から得られた特徴を統合することで，単一の情報源のみを用いる場合よりも，より豊かな表現を獲得することを目的とする．
衣服解析においては，衣服画像に含まれる視覚的特徴と，組成や属性などの数値情報を組み合わせて扱う研究が多く行われている\cite{McAuley15}．画像情報は生地の質感や構造といった視覚的な特徴を含む一方で，数値情報は素材構成などの明示的な属性を表している．これらの情報は互いに補完的であり，両者を統合することで，衣服の特徴をより多面的に捉えることが可能となる．
図\ref{fig:2.1}に，衣服画像および組成情報を統合するマルチモーダル学習の概念図を示す．図に示すように，衣服画像からは画像特徴が，組成情報からは数値特徴がそれぞれ抽出され，これらの特徴は融合処理によって統合表現へと変換される．この統合表現は，各モダリティに含まれる情報を反映した多次元表現であり，以降の処理において共通の特徴表現として用いられる．
マルチモーダル学習における特徴融合の方法には，単純な結合から，各モダリティの寄与を調整する手法まで，様々なアプローチが存在する\cite{Baltrušaitis}．いずれの手法においても，異なる情報源の特徴を適切に統合することが重要であり，融合された表現の質が後続の推定性能に大きく影響することが報告されている\cite{Ngiam}\cite{Baltrušaitis}\cite{Srivastava}\cite{Ramachandram}．


\begin{figure}[tbh] \begin{center}  \scalebox{0.5}{\includegraphics{2-1.PNG}}  \caption{マルチモーダル学習の概要図}  \label{fig:2.1} \end{center}\end{figure}

\section{レビューが表現する感性とテキスト埋め込み}
Web上のレビューには，製品を使用した際の感想や評価が自然言語として記述されており，着心地に関する感性的な情報が含まれている．このようなテキスト情報を計算機で扱うためには，文の意味を数値ベクトルとして表現する手法が必要となる．

近年，自然言語処理分野では，BERT\cite{Devlin}に代表される事前学習済み言語モデルが広く用いられている．BERTは，文脈を考慮した単語表現を学習することができ，文全体の意味をベクトルとして表現することが可能である．これにより，文同士の意味的な類似度を数値的に評価できる．

レビューをベクトル化することで，異なる着心地表現間の関係性を定量的に扱うことが可能となる．本研究では，レビューを教師データとして用いるために，BERTを用いてレビューをベクトル表現に変換する．これにより，異なる着心地表現間の関係性を定量的に扱うことが可能となり，着心地に関する感性的な情報を連続値のベクトルとして扱うことができる．レビューのような自然言語テキストを意味ベクトルとして表現する手法は，文同士の意味的な類似度を数値的に扱うために広く用いられている\cite{Devlin}\cite{Reimers}\cite{Cer}．






\chapter{生地拡大画像と組成情報によるマルチモーダル着心地推定モデル}

\section{レビューを教師信号に用いるマルチモーダル着心地推定}

本研究では，衣服の生地拡大画像および組成情報を入力とし，衣服の着心地を表現するベクトルを推定する，着心地推定モデルを提案する．着心地とは，柔らかさや肌触り，厚みなどの，主観的な感覚を含む概念である．

既存の衣服画像解析に関する研究では，衣服画像を用いたカテゴリ分類や素材認識，質感推定などが主に行われてきた．しかし，これらの研究は主に視覚的属性を対象としており，主観的な着心地に着目した研究は十分に行われてこなかった．

また，レビューを用いた感性分析では，テキストを直接モデルの入力として扱う手法が多い．そのため，レビューが存在しない商品には適用できないという課題がある\cite{McAuley15}\cite{He}\cite{Ni}\cite{McAuley13}\cite{Zheng}．
これに対し，本研究ではレビューをモデルの入力として用いない．代わりに，レビューから生成した文埋め込みベクトルを，着心地を表現する教師データとして利用する．このように，人手で明示的に付与されたラベルではなく，レビューなどの弱い教師信号を用いて学習を行う手法は，既存研究においても有効性が示されている\cite{Zhang}．これにより，学習時にはレビュー情報を活用しつつ，推論時には生地拡大画像と組成情報のみから，着心地推定を行うことが可能となる．この点が，本研究と既存研究との大きな違いである．

以下では，提案手法の全体構成について述べる．図\ref{fig:3.1}に，本研究で提案する着心地推定手法の概要を示す．まず，Web上のショッピングサイトから衣服に関するデータを収集し，生地拡大画像，組成情報，およびレビューを抽出する．生地拡大画像および組成情報は，着心地推定モデルへの入力として用いられる．一方，レビューは文埋め込みモデルを用いてベクトル表現に変換され，教師ベクトルとして利用される．生地拡大画像からは画像特徴が抽出され，組成情報からは素材構成に基づく特徴が抽出される．これらの特徴を統合したマルチモーダルモデルにより，衣服の着心地を表現する推定着心地ベクトルが出力される．モデルの学習では，推定着心地ベクトルと教師ベクトルの類似度を最大化するように最適化を行う．学習後のモデルに対しては，生地拡大画像および組成情報のみを入力とし，推定された着心地ベクトルを得る．得られた着心地ベクトルは，コサイン類似度による評価や近傍分析に用いられ，着心地表現の妥当性を検証する. 

\begin{figure}[tbh] \begin{center}  \scalebox{0.5}{\includegraphics{3-1.PNG}}  \caption{手法概要図}  \label{fig:3.1} \end{center}\end{figure}

\section{生地拡大画の像特徴抽出}

生地拡大画像から着心地に関係する視覚的特徴を抽出するために，本研究ではCNNのエンコーダ部分を用いる．生地拡大画像には，繊維の配置や織り方，表面の凹凸など，局所的な構造が多く含まれており，これらの情報は着心地と強く関連している．CNNは，画像中の局所領域に対して畳み込み演算を行うことで，このような局所構造を段階的に捉えることができるため，生地拡大画像の特徴抽出に適している．また，CNNは，画像内の特徴の位置に対して一定の不変性を持つ．これにより，生地拡大画像において撮影位置やスケールに多少のばらつきが存在する場合でも，安定した特徴表現を得ることが可能である．

入力として用いる生地拡大画像は，Web上のショッピングサイトに掲載されている商品画像のうち，生地部分を拡大して撮影された画像である．各画像は，RGBの3チャネルからなる画素配列として扱い，リサイズおよび中心クロップによる前処理を行った後，画素値を正規化してモデルに入力する．前処理後の画像は，色チャネル数×高さ×幅からなるテンソルとして表現される．これらの画像を CNNエンコーダに入力することで，畳み込み層およびプーリング層を通じて，生地の模様や繊維構造を表す特徴が抽出される．本研究で用いるCNNエンコーダは，画像分類モデルにおける特徴抽出器として設計された構成を基にしており，最終的な分類層は用いない．エンコーダの出力は，全結合層を介して512次元の画像特徴ベクトルへと変換され，後段のマルチモーダルモデルへの入力として用いられる．このように，生地拡大画像をCNNエンコーダによって512次元の特徴ベクトルに変換することで，生地の見た目に関する情報を数値的に表現し，組成情報と統合可能な形式で扱うことができる．


\section{組成情報のベクトル表現}

衣服の着心地は，生地の見た目だけでなく，素材の種類や混率といった物性情報にも大きく依存する．例えば，同じ外観を持つ生地であっても，綿やウール，ポリエステルなどの素材構成が異なる場合，着心地に差が生じることがある．このような情報は，生地拡大画像のみからは十分に推定できない．

そこで本研究では，衣服の組成情報を数値ベクトルとして表現し，生地拡大画像から抽出した特徴と併せてモデルに入力する．組成情報は，Web上のショッピングサイトに記載されている素材の種類および各素材の割合から構成される．
本研究では，組成情報を固定次元の数値ベクトルとして表現し，これを多層パーセプトロン（MLP）に入力する．まず，商品データに出現する素材名の集合をあらかじめ定義し，各素材をベクトルの各次元に対応付ける．各商品の組成表示に対して，対応する次元に混率を代入し，それ以外の次元は0とすることで素材混率ベクトルを構成する．例えば，素材集合を（綿，麻，ポリエステル）とした場合，綿35\%，ポリエステル65\%の生地は（0.35，0.00，0.65）として表現される．

商品によっては，表地および裏地など複数の部位に対して組成表示が与えられる場合がある．本研究では，表地および裏地それぞれについて素材混率ベクトルを構成し両者を加算した後，全成分の総和が1となるよう正規化を行うことで，1つの組成ベクトルへ統合する．これにより，衣服全体としての素材構成比を反映した表現を得る．
構成した組成ベクトルはMLPに入力され，Layer Normalization による正規化と線形変換を経た後，GELU を活性化関数とする非線形変換が施される．さらに，ドロップアウトによる正則化を行い，入力次元と出力次元が一致する場合には残差接続を用いることで，特徴表現の安定化を図っている．MLPの出力は組成特徴ベクトルとして扱い，生地拡大画像から得られた画像特徴ベクトルと統合される．本研究では，組成特徴ベクトルの次元数を512次元とした．


\section{マルチモーダルモデルの構成}

本研究では，生地拡大画像から抽出した画像特徴と，組成情報から抽出した組成特徴を統合し，着心地を表現する特徴ベクトルを推定するマルチモーダルモデルを構築する．
図\ref{fig:3.2}に，本研究で用いるマルチモーダルモデルの構成を示す．生地拡大画像はCNNエンコーダに入力され，画像の質感や表面構造に関する特徴が抽出される．一方，組成情報はMLPに入力され，素材構成に基づく特徴が抽出される．これら2種類の特徴は，それぞれ異なる情報源に基づく表現であり，着心地推定において相補的な役割を果たす．

本研究では，画像特徴と組成特徴を統合する方法として，ゲート機構を用いた特徴融合手法\cite{Arevalo}を採用する．ゲート付きFusionでは，画像特徴および組成特徴に対して重み付けを行い，それぞれの寄与度を調整しながら統合を行う\cite{Baltrušaitis}\cite{Arevalo}．これにより，生地の外観が着心地に強く影響する場合と，素材構成がより重要となる場合の双方に柔軟に対応することが可能となる．画像特徴ベクトル$\mathbf{v}_{\text{img}}$と組成特徴ベクトル$\mathbf{v}_{\text{comp}}$は，ゲート機構を用いて次式により統合される．
$[\cdot;\cdot]$は2つの特徴ベクトルの結合を表し，
$W_g$および$\mathbf{b}_g$ はゲート計算に用いられる学習可能な重みおよびバイアスである．
$\sigma(\cdot)$はシグモイド関数を表し，
$\odot$は要素ごとの積を表す．
$\mathbf{g} \in (0,1)^{512}$は各次元における寄与度を制御するゲートベクトルであり，
$z \in {\cal R}^{512}$は
画像特徴と組成特徴を統合した融合特徴ベクトルである．

\begin{center}

$\mathbf{z}
= \mathbf{g} \odot \mathbf{v}_{\mathrm{img}}
+ (1-\mathbf{g}) \odot \mathbf{v}_{\mathrm{comp}},
\qquad$
$\mathbf{g}
= \sigma\!\left( W_g \,[\mathbf{v}_{\mathrm{img}};\mathbf{v}_{\mathrm{comp}}] + \mathbf{b}_g \right)$

\end{center}

統合された特徴は，全結合層を通じて次元変換され，衣服の着心地を表現する10次元の着心地ベクトルとして出力される．この着心地ベクトルは，着心地を意味空間上の点として表現したものであり，ベクトル間の距離や類似度によって着心地の近さを評価することができる． 

\begin{figure}[tbh] \begin{center}  \scalebox{0.5}{\includegraphics{3-2.PNG}}  \caption{マルチモーダルモデル構成図}  \label{fig:3.2} \end{center}\end{figure}

\section{マルチモーダルモデルの学習}

本研究では，図\ref{fig:3.3}に示すように，pretrainフェーズとfine-tuningフェーズからなる二段階の学習方法を採用する．両フェーズでは同一のマルチモーダルモデル構造を用い， pretrainフェーズで学習したモデルの重みをfine-tuningフェーズに引き継いで学習を行う．

学習に用いる教師データは，レビュー文を文埋め込みモデルを用いて固定次元のベクトルに変換することで生成される．具体的には，日本語事前学習済みのSentence-BERT\cite{Reimers}系モデルを基盤とし，レビュー文をTransformerに入力することで，各トークンに対応する文脈化埋め込みを取得する．得られたトークン埋め込みは平均プーリングによって集約され，文全体を表す文ベクトルへと変換される．さらに，この文ベクトルに対して全結合層を適用し，最終的に10次元の連続値ベクトルとして出力する．この10次元ベクトルを，本研究における着心地表現の教師データとして用いる．

本研究では，着心地ベクトルの次元数を10次元に設定する．これは，着心地を表現する意味空間を過度に高次元化せず，コンパクトな空間として扱うためである．次元数が大きすぎる場合，着心地表現が冗長となり，ベクトル間の距離や類似度に基づく解釈が困難になる可能性がある．
また，出力ベクトルの次元数を抑えることで，モデルの自由度を制限し，学習時の過適合を防ぐ効果が期待できる．特に，fine-tuningフェーズでは教師データ数が限られているため，低次元の着心地空間を用いることで，学習の安定性を向上させることができる.

pretrainフェーズでは，Web上のショッピングサイトから収集したレビュー全文を用いて学習を行う．レビュー全文を文埋め込みモデルで変換することで，大規模な教師ベクトル集合を自動的に生成する．このフェーズでは，生地拡大画像および組成情報を入力とするマルチモーダルモデルの出力である推定着心地ベクトルと，レビュー全文から生成した教師ベクトルとの間のコサイン類似度を最大化するように学習を行う．レビュー全文には着心地以外の情報も含まれるため，教師ベクトルには一定のノイズが含まれるが，大規模データを用いることで，着心地に関する一般的な対応関係を学習することが可能となる．

fine-tuningフェーズでは，レビューのうち着心地に関する記述のみを手動で抽出したデータを用いて再学習を行う．抽出された着心地文を文埋め込みモデルで変換することで，高品質な教師ベクトルを生成する．このフェーズでは，pretrainフェーズで学習したモデルの重みを初期値として用い，推定着心地ベクトルと手動クリーニング済み教師ベクトルとのコサイン類似度を最大化するように追加学習を行う．データ数はpretrainフェーズと比較して少ないが，教師データの品質が高いため，着心地表現に特化した精度向上が期待できる．

このように，大規模データを用いたpretrainと，小規模かつ高品質なデータを用いたfine-tuningを組み合わせることで，着心地に関する汎用的な知識と，着心地表現に特化した精度の両立を実現する． 
本研究では，推定着心地ベクトルと教師ベクトルのコサイン類似度を最大化することを目的として学習を行う．


\begin{figure}[tbh] \begin{center}  \scalebox{0.5}{\includegraphics{3-3.PNG}}  \caption{学習フロー図}  \label{fig:3.3} \end{center}\end{figure}

\chapter{データセットと教師データ構築}
\section{着心地推定タスクの定義}
本研究では，生地拡大画像および組成情報を入力とし，衣服の着心地を表現するベクトルを推定するタスクを着心地推定タスクとして定義する．

入力は，Web上のショッピングサイトに掲載されている衣服の生地拡大画像および組成情報である． 出力は，衣服の着心地を意味空間上の点として表現した10次元の着心地ベクトルである．

本タスクでは，レビューを直接入力として用いない．レビューは，着心地を表現する教師ベクトルを生成するためにのみ利用される．これにより，学習時にレビューを利用しつつ，推論時には生地拡大画像と組成情報のみから着心地推定を行う設定とする．レビューを教師信号として利用する設定は，推薦や感性分析の分野においても採用されており，大規模データを用いた表現学習の有効性が報告されている\cite{Ni}\cite{McAuley13}．

このようなタスク定義により，レビューが存在しない商品に対しても着心地推定が可能となる点が，本研究の特徴である．

このタスク設定は，衣服購入前に着心地を事前に知りたいという実利用上の要求に基づいている．一般的なショッピングサイトでは，生地拡大画像や組成情報は多くの商品で提供されている一方，レビューは商品によっては存在しない場合がある．そのため，推論時にレビューを必要とするタスク設定では，適用可能な商品が限定されてしまう．

また，着心地は数値的な正解値を一意に定義することが困難な主観的概念であるため，本研究では着心地をカテゴリやスコアとして予測するのではなく，意味空間上のベクトルとして表現するタスクを採用した．ベクトル表現を用いることで，着心地の近さや類似性を連続的に扱うことができ，類似商品の検索や比較といった応用にも対応しやすくなる．

このように，本研究のタスク定義は，実世界のデータ可用性と着心地という概念の性質の双方を考慮したものであり，実用性と拡張性を両立することを目的としている．


\section{データセットの概要}

本研究では，第4.1節で定義した着心地推定タスクを検証するために，インターネット上のショッピングサイトに掲載されている衣服商品に関する既存のデータセットを用いる．本節では，使用するデータセットの概要と，本研究で利用する情報の種類について述べる．

本研究の実装例として，楽天が提供する衣服商品データセット \cite{Rakuten} を用いる．本データセットには，楽天市場に掲載されている衣服商品に関する情報として，商品名，商品説明文，商品画像，および購入者によるレビュー文などが含まれている．楽天市場は，多数の衣服商品が掲載されており，衣服の外観や素材構成，着用後の感想に関する情報が比較的充実しているため，着心地推定タスクの検証に適したデータセットである．なお，本研究で提案する着心地推定手法は特定のショッピングサイトに依存するものではなく，同様の情報が提供されている他のショッピングサイトにも適用可能である．

本研究では，これらの情報のうち，生地拡大画像，組成情報，およびレビュー文を利用する．生地拡大画像および組成情報は，着心地推定モデルへの入力として用いられる．レビュー文は，購入者による着用後の感想を含むテキスト情報であり，学習および評価における教師ベクトルを生成する目的のみに使用し，推論時のモデル入力には含めない．


\section{データ抽出条件と前処理}
本研究では，第4.2節で述べた衣服商品データセットの中から，着心地推定タスクに適したデータセットを構築するために，一定の抽出条件および前処理を適用した．本節では，使用したデータの抽出条件と前処理の方針について述べる．

まず，着心地推定モデルの入力として生地拡大画像および組成情報を用いるため，生地拡大画像が存在し，かつ組成情報が明記されている商品を抽出対象とした．生地拡大画像が存在しない商品や，組成情報が欠落している商品は，モデルへの入力情報が不足するため除外した．

次に，学習時に教師ベクトルを生成するため，レビューが存在する商品を抽出対象とした．レビューが存在しない商品は，教師データを構築できないため，学習データセットから除外した．ただし，本研究のタスク定義では，推論時にレビューを入力として用いることはなく，レビューは学習および評価のための教師情報としてのみ利用される．

抽出条件を満たした商品については，生地拡大画像，組成情報，レビューの対応関係を確認し，商品単位で一貫したデータとして整理した．画像データについては，明らかに生地部分が写っていない画像や，解像度が著しく低い画像を除外した．組成情報については，素材名および混率が数値として取得できない場合を除外対象とした．

レビューについては，前処理として不要な記号の除去や表記揺れの統一を行い，文埋め込みモデルへの入力に適した形式に整形した．また，商品説明文や配送に関する記述など，着心地についての言及がないレビューは除外した．

以上の抽出条件および前処理により，生地拡大画像，組成情報，レビューが対応付けられたデータセットを構築した．最終的に，2700件の商品データを使用した．次節では，このデータセットを用いて教師データを自動的に構築する方法について述べる．


\section{教師データの自動構築}
本研究では，第4.3節で構築したデータセットを用いて，着心地推定モデルの学習に用いる教師データを自動的に構築した．本節では，レビューを用いた教師ベクトルの生成方法について述べる．

教師データは，各商品に付随するレビューを文埋め込みモデルによってベクトル表現に変換することで生成する．文埋め込みモデルは，入力された文の意味情報を10次元のベクトルとして表現するモデルであり，レビューに含まれる着心地に関する意味的特徴を捉えることが可能である．

本研究では，各商品に対して複数のレビューが存在する場合，すべてのレビューを一様に用いるのではなく，各商品を代表するレビューを選択して教師データとして利用した．具体的には，各レビューに対してTF-IDFスコアを算出し，その値が高い単語を含む文を当該商品の代表的なレビューとして選択した．

レビューには，着心地に直接関係しない情報が含まれる場合がある．例えば，配送や価格，サイズ感に関する記述などは，着心地推定に対してノイズとなる可能性がある．しかし，本節で構築する教師データはpretrainフェーズにおける大規模学習を目的としているため，一定のノイズを含む教師データであっても，大量のデータを用いることで着心地に関する一般的な対応関係を学習できると考えられる．

このように，レビュー全文を用いて教師データを自動構築することで，人手によるアノテーションを必要とせず，大規模な学習用データセットを構築することが可能となる． 

\section{データ前処理パイプライン}

本節では，図\ref{fig:5.1}に示すデータの前処理パイプラインについて説明する．本研究では，商品データおよびレビューデータに対して段階的な抽出および判定処理を行い，4章で述べた流れに沿って生地拡大画像，組成情報，およびレビューが対応付けられたデータセットを構築した．
まず，商品データ全体からジャンルIDを用いて衣服商品を抽出した．これにより，家電や食品など本研究の対象外となる商品を除外し，衣服に関する商品データのみを対象とした．抽出された衣服商品データに対して，各商品に含まれる画像を入力として，生地拡大画像であるか否かを判定する処理を行った．この判定には，事前に学習したCNNを用いた．具体的には，500枚の画像を目視で確認し，生地部分が明確に写っているかどうかを基準として手動でラベル付けを行った．その後，これらのラベル付き画像を用いて学習したモデルにより，生地部分が明確に写っている画像を含む商品を抽出した．
次に，生地拡大画像を含む商品について，商品説明文から組成情報を抽出した．組成情報が取得可能な商品については，各素材の混率を数値ベクトルとして整理し，組成情報を持つ商品データとして保持した．
一方，レビューデータについては，全レビューから商品URLを参照することで衣服商品に対応するレビューを抽出した．抽出された衣服レビューに対して，レビューの有無を基準として商品を整理し，少なくとも一件以上のレビューを持つ商品を対象とした．さらに，レビュー中に含まれる着心地に関連する語を用いた単語検索を行い，着心地に言及しているレビューを抽出した．
最後に，商品データ側で抽出した組成情報を持つ商品と，レビューデータ側で抽出した着心地に言及するレビューを，商品URLをキーとして対応付けた．これにより，生地拡大画像，組成情報，および着心地に関するレビューが揃った商品データを構築し，以降の学習および評価に用いるデータセットを得た．実際に抽出されたデータ例を図\ref{fig:5.2}と表\ref{tab:5.1}\ref{tab:5.2}に示す．表の番号は，図中に示された数字と対応する．

\begin{figure}[tbh] \begin{center}  \scalebox{0.5}{\includegraphics{5-1.PNG}}  \caption{データ処理パイプライン図}  \label{fig:5.1} \end{center}\end{figure}

\begin{figure}[tbh] \begin{center}  \scalebox{0.5}{\includegraphics{5-2.PNG}}  \caption{生地拡大画像の例}  \label{fig:5.2} \end{center}\end{figure}

\begin{landscape}
    
\begin{table}[tbh]
\begin{center}
\caption{組成情報の例}
\vspace{7mm}
\label{tab:5.1}
\begin{tabular}{|c|c|c|c|c|c|c|c|c|c|c|c|}
\hline
番号 & Acrylic & Cashmere & Cotton & Cupro & Linen & Nylon & Polyester & Polyurethane & Rayon & Tencel & Wool  \\ \hline
1  & 0       & 0        & 60     & 0     & 0     & 0     & 38        & 2            & 0     & 0      & 0     \\ \hline
2  & 0       & 0        & 0      & 0     & 0     & 0     & 100       & 0            & 0     & 0      & 0     \\ \hline
3  & 0       & 0        & 0      & 0     & 0     & 0     & 95        & 5            & 0     & 0      & 0     \\ \hline
4  & 0       & 0        & 0      & 0     & 0     & 0     & 0         & 100          & 0     & 0      & 0     \\ \hline
5  & 0       & 0        & 0      & 0     & 0     & 3.45  & 34.48     & 0            & 0     & 0      & 62.07 \\ \hline
6  & 0       & 0        & 52.5   & 0     & 0     & 0     & 47.5      & 0            & 0     & 0      & 0     \\ \hline
7  & 0       & 0        & 0      & 0     & 100   & 0     & 0         & 0            & 0     & 0      & 0     \\ \hline
8  & 0       & 0        & 0      & 0     & 0     & 10    & 50        & 0            & 0     & 0      & 40    \\ \hline
\end{tabular}
\end{center}
\end{table}

\end{landscape}

\begin{table}[tbh]
\begin{center}
\caption{レビューの例}
\vspace{7mm}
\label{tab:5.2}
\begin{tabular}{|l|p{13cm}|}
\hline
番号 & レビュー                                                                                         \\ \hline
1  & ストレッチなので、余裕でゆったりと着れて、生地もしっかりとしていました。                                                         \\ \hline
2  & もこもこであたたかそう。   でも、デザイン、色合い、手触り、暖かさはとても良いです☆。                                                 \\ \hline
3  & 生地が重いかも。                                                                                     \\ \hline
4  & 生地が柔らかく肌触りがとてもいいです   何度洗ってもホコリが出るのが難点です                                                      \\ \hline
5  & 暖かいし軽いと言っています。   ジャケットより一回り大きめを選びましたが、着せてみるとそれ程大きすぎることもなく軽いので良かったです。 大きすぎず、軽く暖かく、大変喜んでおりました。 \\ \hline
6  & 履いてみると、意外と涼しいので気に入ってます。                                                                      \\ \hline
7  & 麻ですが、ゴワゴワせず、質の良さがわかります。                                                                      \\ \hline
8  & 少し重いですがきれいな生地の暖かいジャケットです。                                                                    \\ \hline
\end{tabular}
\end{center}
\end{table}

\section{生地拡大画像判定モデル}
本節では，生地拡大画像を含む商品を抽出するために用いた生地拡大画像判定モデルについて述べる．生地拡大画像は，衣服の表面構造や質感を視覚的に捉えた画像であり，着心地推定モデルにおいて重要な入力情報となる．しかし，商品ページに掲載されている画像には，全体写真や着用写真などが含まれており，生地拡大画像のみを自動的に識別する必要がある．

本研究では，生地拡大画像であるか否かを判別するために，CNNを用いた画像分類モデルを構築した．生地拡大画像の定義として，生地の表面が拡大され，織りや質感が明確に視認できる画像を正例とし，それ以外の画像を負例とした．

学習用データとして，商品ページから収集した画像の中から500枚を抽出し，目視によって生地拡大画像であるかどうかを判別した．この判別結果に基づいて，各画像に対して手動でラベル付けを行い，生地拡大画像判定モデルの学習データとした．ラベル付けにあたっては，生地部分が画面の大部分を占めており，素材の質感が明確に確認できるかどうかを基準とした．
構築したCNNは，入力画像から特徴を抽出し，生地拡大画像である確率を出力する分類モデルである．学習後のモデルを用いて，衣服商品に含まれる画像を判定し，生地拡大画像を含む商品を抽出した．本研究では，この判定結果を用いて，生地拡大画像を持つ商品データのみを以降の処理対象とした．

本モデルは，生地拡大画像を自動的に抽出するための前処理として利用しており，本研究の主目的である着心地推定モデルの学習には直接用いていない．次節では，この自動構築された教師データの一部を手動で精製し，fine-tuningフェーズに用いる高品質な教師データを構築する方法について述べる．

\section{手動クリーニングによる教師データの精製}
第4.4節で構築した教師データは，着心地についてのレビュー全文を用いて自動生成されたものであり，大規模な学習に適している一方で，着心地と直接関係しない情報を含む場合がある．そのため，本研究ではfine-tuningフェーズに用いる高品質な教師データを構築する目的で，レビューに対する手動クリーニングを行った．

手動クリーニングでは，レビューの中から着心地に関する記述のみを抽出し，教師データとして利用した．具体的には，柔らかさや肌触り，厚み，着用時の快適さなど，着心地を直接表現している文を対象とし，配送や価格，サイズ感，使用シーンに関する記述など，着心地と直接関係しない文は除外した．1文に着心地の記述とそれ以外の記述が混在している場合は，着心地の記述のみの文に修正した．本研究では，700件のデータに手動クリーニングを行った．手動クリーニング後の文が極端に短いデータを除外したため，400件の手動クリーニング済みデータが得られた．

このような手動クリーニングを行うことで，着心地に関する意味情報が明確に反映された教師ベクトルを生成することが可能となる．抽出された着心地文は，第4.4節と同様に文埋め込みモデルを用いてベクトル表現に変換し，fine-tuningフェーズにおける教師データとして利用した．

手動クリーニングによって得られる教師データの数は，自動構築された教師データと比較して少数となる．しかし，教師データの品質が高いため，pretrainフェーズで学習したモデルを着心地推定タスクに特化させる目的において有効である．本研究では，大規模な自動構築教師データによる pretrainと，小規模かつ高品質な手動クリーニング教師データによるfine-tuningを組み合わせることで，学習の安定性と精度の両立を図った．


\section{データセット統計}
本節では，第4.2節から第4.7節までの手順により構築したデータセットの統計情報について述べる．本研究では，pretrainフェーズおよびfine-tuningフェーズにおいて異なる性質の教師データを用いるため，それぞれのデータ規模および内訳を整理する．

pretrainフェーズでは，第4.4節で述べた方法により，レビュー全文を用いて自動構築した教師データを利用した．このデータセットは，生地拡大画像，組成情報，およびレビューが対応付けられた商品から構成されている．各商品に対して複数のレビューが存在する場合，すべてのレビューを一様に用いるのではなく，各商品を代表するレビューを選択して教師データとして利用した. 

この代表文選択は，pretrainフェーズおよびfine-tuningフェーズの両方において同様に適用した．この処理により，冗長な表現や内容の重複を抑えつつ，各商品の特徴をよく表すレビューを教師データとして用いることが可能となる．その結果，教師ベクトル集合の規模を適切に制御しながら，着心地に関する意味情報を効率的に学習に反映させている．fine-tuningフェーズでは，第4.5節で述べた手動クリーニングによって抽出した着心地文を用いた教師データを利用した．このデータセットは，pretrainフェーズで用いたデータセットの一部から構築されており，教師データ数は少数であるが，着心地に関する意味情報が明確に反映された高品質な教師データで構成されている．
本研究では，pretrain用データセットおよびfine-tuning用データセットとは別に，評価に用いるデータセットを構築した．評価用データセットは，学習には使用していない商品データから構成されており，モデルの性能を公平に評価するために用いられる．手動クリーニングされていないデータをpretrain用データセットとし，手動クリーニング済みのデータをfine-tuning用，評価用に分割した．

各データセットに含まれる商品数，レビュー数，および教師ベクトル数の詳細については，表\ref{tab:4.1}に示す．ここで示したデータ数は，生地拡大画像，組成情報，および教師データがすべて揃った商品数を表している．
次章では，本章で構築したデータセットを用いて実際に行った実験設定およびデータ処理の流れについて説明する.

\begin{table}[tbh]
\begin{center}
\caption{データセットの統計}
\vspace{7mm}
\label{tab:4.1}
\begin{tabular}{|l|c|}
\hline
データセット             & データ数 \\ \hline
Pretrain用データセット    & 2000 \\ \hline
Fine-tuning用データセット & 300  \\ \hline
評価用データセット          & 100  \\ \hline
\end{tabular}
\end{center}
\end{table}

\chapter{実験設定}

\section{実験概要}
本章では，第4章で構築したデータセットを用いて実施した実験設定について述べる．本研究の実験では，生地拡大画像および組成情報から衣服の着心地ベクトルを推定するマルチモーダルモデルを学習し，pretrain フェーズおよび fine-tuning フェーズにおける学習挙動を評価した．

実験では，第4章で定義したpretrain用データセットおよびfine-tuning用データセットを用いてモデルの学習を行い，学習にはレビューから生成した教師ベクトルを用いた．また，学習済みモデルの性能評価には，学習に使用していない評価用データセットを用いた．
次節では，生地拡大画像および組成情報を入力とする着心地推定モデルの学習設定について述べる． 

\section{着心地推定モデルの学習設定}
本研究では，生地拡大画像および組成情報を入力とするマルチモーダルモデルを学習するために，pretrainフェーズおよびfine-tuningフェーズの二段階学習を採用した．いずれのフェーズにおいても，モデルは着心地を固定次元のベクトルとして出力し，教師ベクトルとの類似度に基づいて学習を行う．

入力として用いる生地拡大画像は，第3章で述べたCNNエンコーダによって特徴ベクトルに変換される．組成情報については，各素材の混率から構成される組成ベクトルを多層パーセプトロンに入力し，組成の特徴を表すベクトルへと変換する．これら二つの特徴ベクトルは，ゲート付きのFusion機構によって統合され，着心地を表現する低次元の特徴表現が生成される．

モデルの出力は10次元の着心地ベクトルとし，教師データとして用いるレビューから生成した教師ベクトルも同一の次元数に統一した．この設定により，推定ベクトルと教師ベクトルを同一の意味空間上で直接比較することが可能となる．また，次元数を10次元に制限することで，着心地空間をコンパクトに保ち，モデルの過適合を防ぎつつ，学習の安定性を向上させることを目的としている．
学習においては，推定された着心地ベクトルと教師ベクトルとのコサイン類似度を用いた損失関数を採用した．コサイン類似度は，ベクトルの大きさに依存せず意味的な方向の近さを評価できるため，レビューの意味表現を教師とする本研究の設定に適している．モデルは，この類似度が最大化されるようにパラメータを更新することで，生地拡大画像および組成情報から，レビューに基づく着心地表現を再現することを学習する．

pretrainフェーズでは，自動的に構築した大規模な教師データを用いてモデル全体を学習し，生地外観や組成情報と着心地表現との対応関係を広く獲得する．fine-tuningフェーズでは，手動クリーニングによって抽出した高品質な着心地文から生成した教師ベクトルを用い，pretrain フェーズで学習したモデルの重みを初期値として再学習を行う．これにより，着心地に関する意味情報をより明確に反映したモデルへと調整する．

以上の学習設定に基づき，次節では，pretrainフェーズおよびfine-tuningフェーズにおける具体的な実験手順について述べる．

\section{Pretrain および Fine-tuning の実験手順}
本節では，第5.4節で述べた学習設定に基づき，pretrainフェーズおよびfine-tuningフェーズにおける実験手順について述べる．
まず，第4章で構築したデータセットを，pretrain 用データセット，fine-tuning 用データセット，および評価用データセットに分割した．pretrain用データセットには，レビューを用いて自動構築した教師データを含む大規模なデータを用いた．fine-tuning用データセットには，手動クリーニングによって抽出した着心地文から生成した教師データを用いた．評価用データセットは，学習には使用していない商品データから構成されており，学習済みモデルの性能評価のみに用いられる．

pretrainフェーズでは，生地拡大画像および組成情報を入力とし，レビューから生成した教師ベクトルを用いてモデル全体の学習を行った．この段階では，大規模な教師データを用いることで，生地外観や素材構成と着心地表現との一般的な対応関係をモデルに学習させることを目的とした．学習は，設定したエポック数に達するまで，または損失が収束するまで継続した．
次に，pretrainフェーズで学習したモデルの重みを初期値として，fine-tuning フェーズを実施した．fine-tuningフェーズでは，pretrainフェーズと同一のモデル構造および損失関数を用い，手動クリーニングによって得られた高品質な教師データを用いて再学習を行った．これにより，pretrain フェーズで獲得した一般的な特徴表現を保持しつつ，着心地に関する意味情報をより明確に反映させることを目的とした．

各フェーズにおける学習後には，学習済みモデルを用いて評価用データセットに対する推論を行い，生地拡大画像および組成情報から着心地ベクトルを推定した．推定された着心地ベクトルと教師ベクトルとの関係については，次章で定量的および定性的に評価する．
以上の手順により，pretrainフェーズとfine-tuningフェーズを段階的に実施し，着心地推定モデルの学習および評価を行った．


\chapter{実験結果と考察}

\section{評価指標}
本研究では，生地拡大画像および組成情報から推定した着心地ベクトルと，レビューから生成した教師ベクトルとの一致度を評価するために，コサイン類似度を評価指標として用いる. 着心地は主観的かつ連続的な概念であり，ベクトルとして表現されるため，ベクトル間の方向の近さを評価できる指標が適している．

コサイン類似度は，二つのベクトル$v$および$t$に対して，次式で定義される.

\begin{center}
   $cos$\_$sim(v,t)=\dfrac{v\cdot t}{\|v\|\|t\|} $ 
\end{center}

ここで，$v$はモデルによって推定された着心地ベクトルを表し，$t$はレビューから生成した教師ベクトルを表す．コサイン類似度は，ベクトルの大きさに依存せず，意味的な方向の一致度を評価できる指標である．

本研究では，pretrainフェーズおよびfine-tuningフェーズのいずれにおいても，推定された着心地ベクトルと教師ベクトルとのコサイン類似度を用いて評価を行う．コサイン類似度の値が高いほど，推定された着心地表現がレビューに基づく着心地表現と近いことを示す．
以上の評価指標に基づき，次節ではpretrainフェーズおよびfine-tuningフェーズにおける実験結果を示し，提案手法の有効性について検討する．

\section{実験結果}

本節では，第5章で述べた実験設定に基づき実施した実験結果について述べる．本研究では，pretrainフェーズおよびfine-tuningフェーズにおける定量評価としてコサイン類似度を用い，さらに推定された着心地ベクトルの性質を確認するために近傍分析を行った．

pretrainフェーズおよびfine-tuningフェーズにおける平均コサイン類似度の結果を表\ref{tab:6.1} に示す．pretrainフェーズでは平均コサイン類似度が0.48であったのに対し，fine-tuningフェーズでは0.66まで向上した．

この結果から，大規模な自動構築教師データを用いたpretrainフェーズに加えて，手動クリーニングによって構築した高品質な教師データを用いたfine-tuningフェーズを行うことで，推定された着心地ベクトルと教師ベクトルとの一致度が向上していることが確認された．
次に，推定された着心地ベクトルに対して近傍分析を行い，ベクトル空間上における配置の傾向を確認した．近傍分析では，クエリとなるレビューに対し，コサイン類似度に基づき最近傍となるレビューを抽出した．

表\ref{tab:6.2} には，「デザインは気に入っていますが, ちょっと重いです」というクエリレビューに対する近傍分析結果を示す．この例では，生地の重さや厚みに言及したレビューが近傍として抽出されていることが確認された．

表\ref{tab:6.3}  には，「柔らかさとストレッチが効いていて, とても履きやすいです」というクエリレビューに対する近傍分析結果を示す．柔らかさやストレッチ性といった着用感に関する表現を含むレビューが近傍に配置されている様子が確認された．

一方，表\ref{tab:6.4}  には，「生地は暑い時期には不向きです」というクエリレビューに対する近傍分析結果を示す．この例では，クエリレビューと必ずしも同一の着心地表現を含まないレビューが近傍として抽出される例も確認された．

以上の結果から，推定された着心地ベクトルは，一部のクエリに対しては意味的に類似したレビューを近傍として配置する傾向を示す一方で，クエリの内容によっては異なる着心地表現が近傍に配置される場合があることが確認された．これらの結果については，次節で詳細に考察する．

\begin{table}[tbh]
\begin{center}
\caption{コサイン類似度のスコア}
\vspace{7mm}
\label{tab:6.1}
\begin{tabular}{|c|c|}
\hline
学習フェーズ      & 平均コサイン類似度 \\ \hline
Pretrain    & 0.48      \\ \hline
Fine-tuning & 0.66      \\ \hline
\end{tabular}
\end{center}
\end{table}

\begin{table}[tbh]
\begin{center}
\caption{「デザインは気に入っていますが, ちょっと重いです」の近傍データ}
\vspace{7mm}
\label{tab:6.2}
\begin{tabular}{|c|p{9cm}|c|}
\hline
クエリ &
\multicolumn{2}{p{11cm}|}{デザインは気に入っていますが、ちょっと重いです。}
\\ \hline
順位 & 最近傍レビュー & コサイン類似度 \\ \hline
1 & 裏起毛もあり保温性も良く、その割には生地も薄目でゴワゴワしていません。 & 0.93 \\ \hline
2 & 生地が多い分、重いです。 & 0.92 \\ \hline
3 & わりと厚めの生地で裏地も付いているので、暑いかな？布地がたっぷりで裏地もあるので、夏物としては重いかもしれません。生地もずっしり重い感じですが、しっかりしています。 & 0.92 \\ \hline
4 & 生地が重いかも。 & 0.91 \\ \hline
5 & 薄くて暖かそうです。肌触りも良いです。生地もサラサラして着やすいし良いですよ。 & 0.91 \\ \hline
\end{tabular}
\end{center}
\end{table}

\begin{table}[tbh]
\begin{center}
\caption{「柔らかさとストレッチが効いていて、とても履きやすいです。」の近傍データ}
\vspace{7mm}
\label{tab:6.3}
\begin{tabular}{|c|p{9cm}|c|}
\hline
クエリ &
\multicolumn{2}{p{11cm}|}{柔らかさとストレッチが効いていて、とても履きやすいです。}
\\ \hline
順位 & 最近傍レビュー & コサイン類似度 \\ \hline
1 &素材は肌触り心地良い、週末送料無料助かりました。 & 0.95 \\ \hline
2 & ストレッチがきいているのに、1日はいていても膝が出ません。 & 0.95 \\ \hline
3 & とっても軽くて暖かいいい買い物が出来ました。 & 0.94 \\ \hline
4 & パステル調の思った通りの色で、ストレッチが効いていて着やすいです。 サラサラしていて気持ち良くてお気に入りです 色違いも欲しいくらい 柔らかいし、肌の弱い私でも着ててかゆくなりませんでした。 & 0.94 \\ \hline
5 &余所のウルトラストスレッチパンツを持ってますが、 こちらの方がもっとスーパーストレッチでした。 & 0.93 \\ \hline
\end{tabular}
\end{center}
\end{table}

\begin{table}[tbh]
\begin{center}
\caption{「生地は暑い時期には不向きです。」の近傍データ}
\vspace{7mm}
\label{tab:6.4}
\begin{tabular}{|c|p{9cm}|c|}
\hline
クエリ &
\multicolumn{2}{p{11cm}|}{生地は暑い時期には不向きです。}
\\ \hline
順位 & 最近傍レビュー & コサイン類似度 \\ \hline
1 &首元は手にした時から伸びそうで生地が薄く、綿１００％とありましたが、着るとペラペラなのに暑いです。 & 0.95 \\ \hline
2 & 柔らかく動きやすいく暖かいです。 & 0.95 \\ \hline
3 & 着心地もよく暖かいとのことです。 & 0.95 \\ \hline
4 & しっかりした生地で暑いかと思いきや、着るとひんやりして風が吹いたりすると寒いくらいでした。 & 0.94 \\ \hline
5 &生地も柔らかいので、ストレスフリーです。 & 0.94 \\ \hline
\end{tabular}
\end{center}
\end{table}

\section{考察}
本節では，前節で示した実験結果について考察を行う．特に，pretrainフェーズおよびfine-tuningフェーズにおける性能差，ならびに近傍分析における成功例と失敗例の要因について検討する．

まず，定量評価結果について考察する．表\ref{tab:6.1}に示したように，fine-tuning フェーズを行うことで平均コサイン類似度が向上した，この要因として，pretrainフェーズではレビューを用いて自動構築した教師データを使用しており，着心地以外の情報を含むレビューが一定数混在していたことが考えられる．一方，fine-tuningフェーズでは手動クリーニングによって着心地に関する記述のみを抽出した教師データを用いているため，モデルが着心地に関する意味情報をより集中的に学習できたと考えられる．このような，大規模かつノイズを含む教師データによる事前学習と，小規模だが高品質な教師データによる再学習を組み合わせる手法の有効性は，既存研究においても報告されている\cite{Hinton}.

次に，近傍分析における成功例について考察する．表\ref{tab:6.2}および表\ref{tab:6.3} に示した例では，クエリレビューと意味的に類似した着心地表現を含むレビューが近傍として抽出されている．これらのクエリは，「重い」「柔らかい」「ストレッチが効いている」といった，生地の物理的特性や触感に直接対応する表現を含んでいる．このような着心地表現は，生地拡大画像に含まれる視覚的特徴や，組成情報に含まれる素材構成と関連付けやすいため，モデルが適切な着心地ベクトルを推定できたと考えられる．画像から質感や素材に関する情報を抽出できることは，既存の質感認識や素材認識に関する研究においても報告されている\cite{Cimpoi}\cite{Bell}．

一方で，表\ref{tab:6.4}に示した失敗例では，クエリレビューと必ずしも一致しない着心地表現を含むレビューが近傍として抽出されている．「暑い時期には不向き」といった表現は，温熱環境や着用シーン，個人の主観などに依存する側面が大きく，生地拡大画像や組成情報のみから推定することが難しい情報である．このような文脈依存的な感性表現は，マルチモーダル学習においても扱いが難しいことが指摘されている\cite{Baltrušaitis}\cite{Poria16}.この結果は，本研究の手法が生地の視覚的特徴および素材構成に基づく着心地推定には有効である一方で，環境要因や使用状況に強く依存する着心地表現の推定には限界があることを示している．このような文脈依存的な感性や感情の推定は，アフェクティブコンピューティングの分野においても困難であることが指摘されている\cite{Schuller}．

また，本研究の手法には，教師データとしてレビューを用いていることに起因する限界が存在する．本手法では，レビューに明示的に記述されている着心地情報のみを教師信号として学習を行っているため，レビューに書かれていない着心地に関する情報を考慮することができない．例えば，着用者が無意識に感じている快適さや，特定の環境下でのみ顕在化する着心地の違いなどは，レビューとして表現されない場合がある．

このような情報は，教師データとしてモデルに与えられないため，生地拡大画像や組成情報に含まれていたとしても，本研究の枠組みでは十分に反映されない可能性がある．したがって，本手法はレビューに基づいて言語化された着心地表現を推定することには有効である一方で，言語化されていない潜在的な着心地情報の推定には限界を有している．このような主観的・感性的概念を計算機上で定量的に扱うことの難しさは，感性工学やアフェクティブコンピューティング，および経験価値に関する研究分野においても指摘されている\cite{Ramachandram}\cite{Cer}\cite{Poria17}\cite{Schmitt}．

以上の考察から，本研究で提案した手法は，生地の質感や構造に基づく着心地の推定において有効であることが確認された．しかし，温度感覚や季節適性などの要素については，画像および組成情報以外の情報を導入する必要があると考えられる．例えば，商品説明文や使用シーンに関する情報を追加的に用いることで，これらの要因を考慮した着心地推定が可能になると考えられる．



        
\chapter{おわりに}
                
本研究では，生地拡大画像と組成情報を用いて衣服の着心地を推定する手法を提案した．インターネット上のショッピングサイトにおいては，実際に衣服を手に取ることができないという制約があり，購入前に着心地を正確に把握することは困難である．この問題に対し，本研究では視覚情報と数値情報を統合したマルチモーダル学習により，着心地をベクトルとして表現する枠組みを構築した．

提案手法では，生地拡大画像からCNNを用いて視覚的特徴を抽出し，組成情報から多層パーセプトロンを用いて数値的特徴を抽出した．これらの特徴を融合することで，衣服の着心地を10次元のベクトル空間に写像した．教師データとしては，着心地に関するレビューをSentence-BERTによってベクトル化し，大規模データによる事前学習と，高品質データによる再学習を段階的に行った．

実験の結果，再学習を行うことで，推定された着心地ベクトルと教師ベクトルとのコサイン類似度が0.48から0.66に向上し，表現の精度が改善されることを確認した．また，近傍分析により，推定された着心地ベクトルが，レビューに含まれる着心地に関する意味的な類似性を適切に捉えていることを定性的に示した．これらの結果から，提案手法が生地拡大画像および組成情報から着心地を推定する上で有効であることが示唆された．

一方で，本研究にはいくつかの課題も存在する．本研究の手法はレビューを教師データとして用いて学習を行っているため，言語化されていない潜在的な着心地情報については推定が困難である．例えば，個人の体型差や着用時の環境条件，長期間の着用による変化などは，レビューに明示的に記述されない場合が多く，本手法の枠組みでは十分に捉えられない可能性がある．

今後の課題としては，教師データの拡充や，レビュー以外の情報源を取り入れた着心地表現の高度化が挙げられる．例えば，時系列情報やユーザ属性を考慮することで，より個人に適した着心地推定が可能になると考えられる．また，推定された着心地ベクトルを自然言語として生成することで，ユーザにとって理解しやすい形で提示することも今後の展開として考えられる．

以上より，本研究は，生地拡大画像と組成情報を用いたマルチモーダル学習により，衣服の着心地をベクトルとして表現する新たな枠組みを示した．本手法は，オンラインショッピングにおける着心地理解を支援する基盤技術として，今後の発展が期待される．




\chapter*{謝辞}

本稿をまとめる動機を与え、さらに、執筆を支えて
くださった立命館大学情報理工学部データ工学研究室のみなさまに心より感謝し
ます。

本研究では、国立情報学研究所のIDRデータセット提供サービスにより楽天グループ株式会社から提供を受けた「楽天データセット」（https://rit.rakuten.com/data\_release/）を利用しました。

\begin{thebibliography}{99}

% 参考文献は，first author の family name のアルファベット順に並べるのが一般的

\addcontentsline{toc}{chapter}{参考文献}
%
\bibitem{Park}Park， J.， Lennon， S. J.， and Stoel， L.，“On-line product presentation: Effects on mood， perceived risk， and purchase intention，”Psychology \& Marketing， 2005.
%
\bibitem{LiuZ}Liu, Z. et al., “DeepFashion: Powering Robust Clothes Recognition and Retrieval,” CVPR, 2016.
%
\bibitem{Bell}Bell, S. et al., “Material Recognition in the Wild with the Materials in Context Database,” CVPR, 2015.
%
\bibitem{McAuley15} McAuley, J. et al., “Image-based Recommendations on Styles and Substitutes,” ICDM, 2015.
%
\bibitem{He}He, R. et al., “Ups and Downs: Modeling the Visual Evolution of Fashion Trends,” WWW, 2016.
%
\bibitem{Ngiam}Ngiam, J. et al., “Multimodal Deep Learning,” ICML, 2011.
%
\bibitem{Baltrušaitis}Baltrušaitis, T. et al., “Multimodal Machine Learning: A Survey,” TPAMI, 2019.
%
\bibitem{Ni}Ni, J. et al., “Justifying Recommendations using Reviews and Images,” SIGIR, 2019.
%
\bibitem{Norman} Norman， D. A.，“Emotional Design: Why We Love (or Hate) Everyday Things，”Basic Books， 2004.] Norman， D. A.，“Emotional Design: Why We Love (or Hate) Everyday Things，”Basic Books， 2004.
%
\bibitem{McAuley13}McAuley, J. and Leskovec, J., “Hidden Factors and Hidden Topics,” RecSys, 2013.
%
\bibitem{Zheng}Zheng, L. et al., “Joint Learning of Review-Level Sentiment and Aspect,” EMNLP, 2017.
%
\bibitem{LiuQ}Liu, Q. et al., “Deep Learning for Fashion Recommendation: A Survey,”ACM Computing Surveys, 2021.

%
\bibitem{Wu} Wu, S. et al., “Multimodal Fashion Recommendation with Deep Learning,” Information Processing \& Management, 2020.
%
\bibitem{LeCun}LeCun， Y.， Bottou， L.， Bengio， Y.， and Haffner， P.，“Gradient-Based Learning Applied to Document Recognition，” Proceedings of the IEEE， 1998.
%
\bibitem{Cimpoi}Cimpoi, M. et al., “Deep Filter Banks for Texture Recognition,” CVPR, 2015.
%\cite{Ngiam}\cite{Baltrušaitis}}\cite{Srivastava}\cite{Ramachandram}
\bibitem{Sharan}Sharan， L.， Rosenholtz， R.， and Adelson， E. H.，“Accuracy and speed of material categorization in real-world images，”Journal of Vision， 2014.
%
\bibitem{Srivastava}Srivastava, N. and Salakhutdinov, R., “Multimodal Learning with Deep Boltzmann Machines,” NeurIPS, 2012.
%
\bibitem{Ramachandram}Ramachandram, D. and Taylor, G., “Deep Multimodal Learning,” IEEE SPM, 2017.
%
\bibitem{Devlin}Devlin, J. et al., “BERT,” NAACL, 2019.
%
\bibitem{Reimers}Reimers, N. and Gurevych, I., “Sentence-BERT,” EMNLP, 2019.
%
\bibitem{Cer}Cer, D. et al., “Universal Sentence Encoder,” EMNLP, 2018.
%
\bibitem{Zhang}Zhang， Y.， Chen， T.， and Wang， H.，“Learning from Weak Supervision for Product Attribute Extraction，”ACL， 2014.
%
\bibitem{Arevalo}Arevalo，J.，Solorio，T.，Montes-y-Gómez，M.，and González，F.，
“Gated Multimodal Units for Information Fusion，”
ICLR Workshop，2017.
%
\bibitem{Rakuten}See the description in each dataset's item.Use the following only when citing the entire collection.Rakuten Group, Inc. (2014): Rakuten Dataset. Informatics Research Data Repository, National Institute of Informatics. (dataset). https://doi.org/10.32130/idr.2.0
%
\bibitem{Hinton}Hinton， G. et al.，“Distilling the Knowledge in a Neural Network，”NIPS Workshop， 2015.
%
\bibitem{Poria16}Poria， S. et al.，“Context-Dependent Sentiment Analysis，”ACL， 2016.
%
\bibitem{Schuller}Schuller， B. et al.，“Recognising realistic emotions and affect in speech: State of the art and lessons learnt，”Speech Communication， 2011.
%
\bibitem{Poria17}Poria, S. et al., “A Review of Affective Computing,” TKDE, 2017.
%
\bibitem{Schmitt}Schmitt, B., “Experiential Marketing,” Journal of Marketing Management, 1997.
%




\end{thebibliography}
%

\end{document}




